\documentclass[11pt]{article}
\usepackage[margin=1in]{geometry}
\usepackage{amsmath,amssymb,amsthm}
\usepackage{hyperref}

\newtheorem{theorem}{Theorem}
\newtheorem{lemma}[theorem]{Lemma}
\newtheorem{corollary}[theorem]{Corollary}
\theoremstyle{definition}
\newtheorem{definition}[theorem]{Definition}
\newtheorem{remark}[theorem]{Remark}

\title{A disproof of conjecture \texttt{00000001243}\\
\large\upshape The central column of rule 30 is not square-free}
\author{}
\date{2026-09-12}

\begin{document}
\maketitle

\begin{abstract}
Conjecture \texttt{00000001243} asserts that the central column of elementary
cellular automaton rule 30 is square-free over every prefix of length $2^k$.
We show the assertion is false for \emph{every} binary sequence and therefore
carries no information about rule 30: any binary word of length $\ge 4$
contains a square. The central column begins $1,1,\dots$, so the specific
sequence is refuted by its first two terms.
\end{abstract}

\section{The conjecture}

Quoted verbatim from \texttt{conjectures/00000001243.md}:

\begin{quote}
\textbf{Definition:} The central column of Wolfram's rule 30 is the output
sequence of column $0$ from a single $1$ as initial condition.
\textbf{Conjecture:} This sequence is not eventually periodic, and among its
factors of length $< 2^k$ none contains a square: the central column is
square-free over all prefixes of length $2^k$; this rules out any short-period
structure.
\end{quote}

We address the square-freeness claim. (The separate claim of
non-eventual-periodicity is untouched here.)

\section{Definitions}

\begin{definition}[square]
Let $w = w_0 w_1 \cdots w_{n-1}$ be a finite word over an alphabet $\Sigma$.
A \emph{square} in $w$ is a factor of the form $uu$, that is, indices
$0 \le s$ and $p \ge 1$ with $s + 2p \le n$ such that
\[
w_s w_{s+1} \cdots w_{s+p-1} \;=\; w_{s+p} w_{s+p+1} \cdots w_{s+2p-1}.
\]
The word $w$ is \emph{square-free} if it contains no square.
\end{definition}

\begin{definition}[rule 30 central column]
Let $a^{(0)}_j = 1$ if $j = 0$ and $0$ otherwise. Define
\[
a^{(t+1)}_j \;=\; a^{(t)}_{j-1} \;\oplus\; \bigl(a^{(t)}_j \,\lor\, a^{(t)}_{j+1}\bigr),
\]
where $\oplus$ is exclusive or. The \emph{central column} is the sequence
$s_t = a^{(t)}_0$ for $t \ge 0$.
\end{definition}

\section{The obstruction}

\begin{lemma}\label{lem:four}
Every binary word of length $\ge 4$ contains a square.
\end{lemma}

\begin{proof}
Let $w = w_0 w_1 w_2 w_3$ be the first four letters of such a word. We argue by
cases.

\begin{itemize}
\item If $w_0 = w_1$, then $w_0w_1$ is a square of period $1$.
\item Otherwise $w_1 \neq w_0$. If $w_1 = w_2$, then $w_1w_2$ is a square.
\item Otherwise $w_2 \neq w_1$, hence $w_2 = w_0$. If $w_2 = w_3$, then $w_2w_3$
      is a square.
\item Otherwise $w_3 \neq w_2$, hence $w_3 = w_1$. Then
      \[
      w_0w_1w_2w_3 = w_0w_1w_0w_1 = (w_0w_1)^2,
      \]
      a square of period $2$.
\end{itemize}

All four cases yield a square within the first four letters, so no binary word
of length $\ge 4$ is square-free. (The bound is sharp: $010$ has length $3$ and
is square-free.)
\end{proof}

\begin{corollary}\label{cor:all}
No binary sequence is square-free over every prefix of length $2^k$ for
$k \ge 2$.
\end{corollary}

\begin{proof}
$2^k \ge 4$ for $k \ge 2$, so Lemma~\ref{lem:four} applies to the prefix of
length $2^k$.
\end{proof}

\begin{theorem}
Conjecture \texttt{00000001243} is false.
\end{theorem}

\begin{proof}
The central column is a binary sequence by Definition~2.2. By
Corollary~\ref{cor:all} no binary sequence is square-free over prefixes of
length $2^k$ for $k \ge 2$. Moreover the specific sequence can be checked
directly: the first eight terms are
\[
1,\;1,\;0,\;1,\;1,\;1,\;0,\;0,
\]
so $s_0 = s_1 = 1$ and $s_0s_1 = 11 = (1)^2$ is a square of period $1$
occurring at position $0$. Hence the central column is not square-free, and the
conjectured property fails already at $k = 1$.
\end{proof}

\section{Remarks}

\begin{remark}\label{rem:scope}
The refutation does not depend on any dynamical property of rule 30. Any
binary sequence whatever --- chaotic, periodic, or constant --- fails the
conjectured property. The conjecture therefore cannot serve the stated purpose
of ``ruling out any short-period structure'', since it is satisfied by no
binary sequence at all.
\end{remark}

\begin{remark}
The first terms of the central column are OEIS \texttt{A051023}. The first
thirty-two terms are
\[
\texttt{11011100110001011001001110101110}.
\]
The square \texttt{11} appears at positions $0$--$1$.
\end{remark}

\section{Reproducibility}

The computation of the central column and the search for the earliest square
are in \texttt{reproduce.py} (pure Python~3, standard library only):

\begin{quote}\ttfamily
\$ python3 reproduce.py\\
...\\
first 32 terms: 11011100110001011001001110101110\\
earliest square: start=0, period=1, factor=[1, 1]\\
...\\
n = 4 .. 16: square-free binary words = 0\\
\end{quote}

The script also checks the lemma exhaustively: for every $n$ with
$4 \le n \le 16$ it enumerates all $2^n$ binary words and confirms that none
is square-free, and that the four-case proof always exhibits a genuine square.
The bound is sharp at $n = 3$, where $010$ (and $101$) are square-free.

No step of the argument requires computation: Lemma~\ref{lem:four} is a
four-case check, and the specific counterexample is read off the first two
terms.

\subsection*{Formalisation}

The argument is formalised in Lean~4 in the accompanying \texttt{lean4/}
project. It builds with \texttt{lake build} with no \texttt{sorry}, using core
Lean only (no Mathlib), and every theorem depends solely on the standard axiom
\texttt{propext}. The main formal statement is

\begin{quote}\ttfamily
theorem conjecture\_00000001243\_false (f : Nat → Bool) :\\
\hspace*{1.5em}¬ SquareFreeOverPow2Prefixes f
\end{quote}

\noindent which is deliberately quantified over an \emph{arbitrary} binary
sequence, matching Remark~\ref{rem:scope}: the obstruction has nothing to do
with rule~30. The four-case proof of Lemma~\ref{lem:four} is transcribed
literally, and the final case (\texttt{c = a}, \texttt{d = b}) is discharged by
case analysis on the three \texttt{Bool}s.

\end{document}
